\section{Introduction}

Let $h>0$ and let $(v_x)_{x\in\Z^2}$ be independent random variables,
each uniform on $[-h,h]$. We consider the bounded self-adjoint operator
\begin{equation}\label{setup:model}
 (H_v u)(x)=\sum_{\abs{y-x}_1=1}u(y)+v_xu(x)
 \qquad\text{on }\ell^2(\Z^2).
\end{equation}
Thus $v_x=h\omega_x$ in the usual notation with $\omega_x$ uniform on
$[-1,1]$. Pure-point spectral type means that the closed span of the
eigenvectors is the entire Hilbert space; it does not mean that every
point of the spectrum is an eigenvalue.

The planar spectral localization problem asks whether this pure-point
conclusion holds throughout the spectrum for every nonzero width of the
uniform distribution \cite[Example~2.1 and Problem~2]{Simon2000}.
We prove that it does.

\begin{theorem}\label{main:theorem}
For each fixed $h>0$, almost surely $H_v$ has pure-point spectral type,
admits a complete orthonormal eigenbasis, and has spectrum
\[
 \sigma(H_v)=[-4-h,4+h].
\]
Its eigenvalues are dense in this interval. The probability-one event
in this statement is allowed to depend on $h$.
\end{theorem}

The conclusion includes arbitrarily weak nonzero disorder and the entire
spectrum. It resolves the pure-point assertion of the two-dimensional
Anderson localization conjecture for a uniform single-site distribution,
as formulated in \cite[Problem~2]{Simon2000}. An arbitrary nondegenerate
uniform interval reduces to the centered interval above by subtracting
its midpoint from the operator. We do not assert one probability-one
event for all disorder strengths, a bound uniform in energy, or a
dynamical-localization theorem.

\subsection*{Context and scope}
Anderson introduced random site energies as a mechanism for suppressing
quantum transport, motivated by spin diffusion and impurity conduction
\cite{Anderson1958}. His discussion includes a flat single-site
distribution and nearest-neighbor hopping, although the original
physical argument is not a modern proof of the spectral assertion
considered here. The conductance-scaling theory of Abrahams, Anderson,
Licciardello and Ramakrishnan predicts the absence of true metallic
behavior in two dimensions
\cite{AbrahamsAndersonLicciardelloRamakrishnan1979}. That physical
prediction is distinct from an almost-sure spectral-type theorem.
Simon's formulation isolates the latter question; his diffusion
problem is listed separately \cite{Simon2000}.

Rigorous multidimensional localization theory developed several
complementary mechanisms.
Fr\"ohlich and Spencer's multiscale analysis gave Green-function decay
and vanishing Kubo diffusion in large-disorder or appropriate
extreme-energy regimes \cite{FrohlichSpencer1983}. Their 1983
conclusions were explicitly distinguished from a complete pure-point
decomposition. Fr\"ohlich, Martinelli, Scoppola and Spencer subsequently
established pure-point spectrum with exponentially decaying
eigenfunctions for independent identically distributed diagonal
potentials with a bounded density in sufficiently large-disorder and suitable tail-energy
regimes \cite{FMSS1985}. Delyon, L\'evy and Souillard also obtained
multidimensional spectral localization in large-disorder or large-energy
regimes \cite{DelyonLevySouillard1985}. The passage from fixed-energy
Green-function information to a complete spectral decomposition is a
substantive step, not an automatic consequence of decay outside an
exceptional set of energies.

Aizenman and Molchanov developed the fractional-moment approach,
obtaining localization at large disorder and at extreme energies under
suitable hypotheses on the conditional single-site distributions,
including absolute continuity and regularity
\cite{AizenmanMolchanov1993}. This method supplies important ancestry for
our use of fractional Green-function moments, but its established
regimes do not by themselves yield the all-disorder, all-energy
assertion above. Spectral-edge results also have a different scope.
Ding and Smart proved localization near an edge for the two-dimensional
Bernoulli model and observed that their proof extends to bounded
nontrivial iid laws
\cite{DingSmart2020}; Hurtado extended edge-localization and
unique-continuation methods to independent, not necessarily identically
distributed uniformly bounded potentials with uniformly positive variance
\cite{Hurtado2026}. The latter
localization theorems remain restricted to a bottom interval, even
though their distributional assumptions are broader.

Weak-disorder transport on long finite scales poses another, compatible
question. Black, Drogin and Hern\'andez obtain quantitative time-averaged
diffusion, resolvent and finite-volume eigenfunction results for
Gaussian iid potentials on disorder-dependent scales
\cite{BlackDroginHernandez2025}.
These weak-coupling limits neither imply nor contradict pure-point
spectral type at every fixed positive disorder for the uniform law.
The present assertion concerns that fixed-disorder spectral type; it
does not assert exponential eigenfunction decay or dynamical localization.

The final spectral step follows the rank-one mechanism of Simon and
Wolff \cite{SimonWolff1986}. Its input is a suitable inverse-square
spectral integral, equivalently a square-summable resolvent column, for
almost every configuration and energy. Averaging over the independent
site value transfers this information to the random spectral measure.
We prove the required averaging and spectral readout explicitly; the
main work is to obtain their fixed-energy input under the true uniform
product law at arbitrary disorder.

\subsection*{The proof mechanism}
Fix $h>0$ and an interior real energy $E$. In \cref{mat:section}, we introduce absorbing
diagonal terms in finite tori and their associated input and output
channels, called ports. Squared scattering amplitudes define a
Dirichlet energy on port data: if $S_{ja}$ is the scattering amplitude
from port $a$ to port $j$ and $f$ is real data on the ports, the energy is
\[
 \mathcal Q_S(f)=\frac12\sum_{a,j}|S_{ja}|^2|f_j-f_a|^2.
\]
A Cauchy-distributed real closing value
would remove a port with a favorable averaged energy comparison. The
bounded physical distribution requires a truncated version and a
relative estimate for its error.

The first main ingredient is a family of rare resonant cells
(\cref{sec:resonant,sec:impedance,sec:transmission}). Their
conditional law is embedded in the original product measure by an
acceptance likelihood; the true potentials remain independent uniforms
before conditioning. The construction gives a weighted boundary
impedance estimate uniform in arbitrary passive exterior data, including
weakly coupled missing boundary directions. It also gives a transmission
lower bound strong enough to make the closing error small relative to
the transmission itself. These estimates, together with an upper
envelope over adaptive interior configurations, are useful beyond a
single prescribed environment.

The second ingredient is a comparison along a hierarchy of increasingly
sparse ports (\cref{sec:hierarchy}). A parent-cell alternative controls the apparent singularity
at the end of each thinning interpolation. Local generic observation,
followed by path estimates, supplies screening at every fixed starting
stage (\cref{sec:screening}). For a fixed starting level $n$, let $d_n$
be the infinite-volume limit of the expected sum of squared port-to-port
distances, weighted by $|S_{ja}|^2$ and divided by twice the torus area.
These quantities measure the scattering energy of affine coordinate
data. The torus side tends to infinity with $n$ fixed to define $d_n$;
only afterwards does the thinning comparison give a finite limit $d$
as $n$ tends to infinity.

The dimension enters decisively in forcing $d=0$ (\cref{sec:vanishing}).
An adaptive modification of affine data first makes a local affine
first-variation trace small in expectation. Separately, a logarithmic voltage profile separates
retained interior ports from an annulus of very weak probes at
arbitrarily small two-dimensional capacity cost. If $d>0$,
an analytic test of the scattering matrix has a nonzero value at the
center of the disk. Support-preserving shifts of the true potentials
and the thinning comparison make its boundary moments arbitrarily small,
contradicting submean. Dimension-free weak integral bounds prevent the
number of ports from spoiling this argument.

Finally, vanishing affine energy is upgraded to a fixed-real-energy
fractional-moment estimate for the true finite-torus Green function
(\cref{sec:bootstrap}).
Buffered probe barriers, conditional estimates under transmission
weights, and resampling of disjoint regions close a contraction uniform
over the allowed masks. The infinite-volume conclusion then follows
from the spectral theorem, Fubini, and rank-one averaging
(\cref{sec:spectral}). The property
to which Fubini is applied belongs to the true operator, so no common
auxiliary construction measurable in energy is needed.

The proof keeps these roles separate: the rare laws prove local
comparisons, fixed-stage screening controls limits, two-dimensional
capacity makes the limit vanish, and spectral averaging converts
fixed-energy information to spectral type.
